% ============================================================================
% Chapter 2 — Classical Arabic Poetry (الشعر الفصيح)
% ============================================================================
\chapter{Classical Arabic Poetry}
\label{ch:classical}

Classical Arabic poetry (\ar{الشعر العربي الفصيح}) is governed by one of the most rigorously codified prosodic systems in literary history. This chapter documents that system in full, as required for AI generation and evaluation.

% ============================================================================
\section{Foundations of Arabic Prosody (\ar{علم العروض})}
\label{sec:classical_prosody_foundations}

Arabic prosody was codified by \textbf{al-Khalil ibn Ahmad al-Farahidi} (718–786 CE) in his treatise \textit{Kitab al-'Arud}. His system describes all valid Arabic poetic meters through combinations of three fundamental prosodic units:

\subsection{The Three Prosodic Building Blocks}
\label{subsec:building_blocks}

\begin{table}[H]
\centering
\caption{Fundamental prosodic units (\ar{مكونات العروض})}
\label{tab:prosodic_units}
\renewcommand{\arraystretch}{1.5}
\begin{tabularx}{\textwidth}{llXll}
\toprule
\textbf{Unit} & \textbf{Arabic} & \textbf{Definition} & \textbf{Pattern} & \textbf{Example} \\
\midrule
Sabab Khafif & \ar{سبب خفيف} & Moving + quiescent & 1\,0 & \ar{مَنْ} \\
Sabab Thaqil & \ar{سبب ثقيل} & Moving + moving & 1\,1 & \ar{مَعَ} (rare) \\
Watad Majmu' & \ar{وتد مجموع} & Moving + moving + quiescent & 1\,1\,0 & \ar{فَعُو} \\
Watad Mafruq & \ar{وتد مفروق} & Moving + quiescent + moving & 1\,0\,1 & \ar{مَفْعُ} \\
Fasila Sughra & \ar{فاصلة صغرى} & Three moving + quiescent & 1\,1\,1\,0 & \ar{فَعِلُنْ} \\
Fasila Kubra & \ar{فاصلة كبرى} & Four moving + quiescent & 1\,1\,1\,1\,0 & \ar{فَعَلَتُنْ} \\
\bottomrule
\end{tabularx}
\end{table}

\noindent\textbf{Binary notation convention:} Throughout this document, \textbf{1} denotes a moving (\ar{متحرك}) position (a consonant bearing any short vowel), and \textbf{0} denotes a quiescent (\ar{ساكن}) position (a consonant bearing sukun, or a long vowel \ar{ا/و/ي}). Long vowels are inherently quiescent in the scan.

\subsection{The Eight Tafeelat}
\label{subsec:eight_tafeelat}

Al-Khalil derived eight canonical feet (\ar{تفعيلات}) from these building blocks. All sixteen meters are constructed from these eight feet:

\begin{table}[H]
\centering
\caption{The eight canonical feet (\ar{التفعيلات الثمانية})}
\label{tab:eight_tafeelat}
\renewcommand{\arraystretch}{1.6}
\begin{tabularx}{\textwidth}{lllXl}
\toprule
\textbf{Name (Arabic)} & \textbf{Name (Transcription)} & \textbf{Pattern} & \textbf{Components} & \textbf{Used In} \\
\midrule
\ar{فَعُولُنْ} & Fa'ulun & 1\,1\,0\,1\,0 & Watad majmu' + sabab khafif & Taweel, Mutakareb \\
\ar{مَفَاعِيلُنْ} & Mafa'ilun & 1\,1\,0\,1\,1\,0 & Watad majmu' + sabab thaqil + sabab khafif & Taweel, Hazaj \\
\ar{مَفَاعَلَتُنْ} & Mafa'alatun & 1\,1\,0\,1\,1\,1\,0 & Watad majmu' + fasila sughra + ... & Wafir \\
\ar{فَاعِلُنْ} & Fa'ilun & 1\,0\,1\,1\,0 & Sabab khafif + watad majmu' & Baseet, Taweel \\
\ar{فَاعِلَاتُنْ} & Fa'ilatun & 1\,0\,1\,1\,0\,1\,0 & Sabab khafif + watad majmu' + sabab khafif & Ramal, Khafeef \\
\ar{مُسْتَفْعِلُنْ} & Mustaf'ilun & 1\,0\,1\,0\,1\,1\,0 & Sabab khafif + sabab khafif + watad majmu' & Baseet, Rajaz, Saree \\
\ar{مُتَفَاعِلُنْ} & Mutafa'ilun & 1\,1\,1\,0\,1\,1\,0 & Fasila sughra + watad majmu' & Kamel \\
\ar{مُسْتَفْعِ لُنْ} & Mustaf'i lun & 1\,0\,1\,0\,1\,0 & Watad mafruq + sabab khafif & Khafeef, Mujtath \\
\bottomrule
\end{tabularx}
\end{table}

\subsection{The Five Prosodic Circles}
\label{subsec:five_circles}

Al-Khalil organised the meters into five abstract \textbf{prosodic circles} (\ar{الدوائر العروضية الخمس}). Meters in the same circle share the same cyclic root and differ only in the starting point of the reading:

\begin{table}[H]
\centering
\caption{The five Khalilian prosodic circles}
\label{tab:five_circles}
\renewcommand{\arraystretch}{1.4}
\begin{tabularx}{\textwidth}{llXl}
\toprule
\textbf{Circle} & \textbf{Arabic Name} & \textbf{Meters} & \textbf{Dominant Foot Type} \\
\midrule
1 & \ar{دائرة المختلف} & Taweel, Madeed, Baseet & Mixed (watad + sababaan) \\
2 & \ar{دائرة المؤتلف} & Wafir, Kamel & Fasila-based \\
3 & \ar{دائرة المجتلب} & Hazaj, Rajaz, Ramal & Repeated sabab-watad \\
4 & \ar{دائرة المشتبه} & Saree, Munsareh, Khafeef, Mudhare, Muqtadheb, Mujtath & Contains watad mafruq \\
5 & \ar{دائرة المتفق} & Mutakareb, Mutadarak & Pure trisyllabic feet \\
\bottomrule
\end{tabularx}
\end{table}

% ============================================================================
\section{The Sixteen Khalilian Meters}
\label{sec:sixteen_meters}

Each meter entry below provides: (1) the canonical tafeelat for both hemistichs; (2) the binary pattern; (3) the main zehafs and their effect; (4) the valid aruz/dharb combinations; (5) typical themes; and (6) a famous example verse with full diacritics.

% -----------------------------------------------------------------------
\subsection{\ar{البحر الطويل} — Bahr al-Taweel (``The Long'')}
\label{subsec:taweel}

\begin{table}[H]
\centering
\caption{Al-Taweel structure}
\renewcommand{\arraystretch}{1.4}
\begin{tabularx}{\textwidth}{lX}
\toprule
\textbf{Hemistich (Sadr/Ajuz)} & \ar{فَعُولُنْ مَفَاعِيلُنْ فَعُولُنْ مَفَاعِيلُنْ} \\
\textbf{Pattern} & 1\,1\,0\,1\,0 \; 1\,1\,0\,1\,1\,0 \; 1\,1\,0\,1\,0 \; 1\,1\,0\,1\,1\,0 \\
\textbf{Circle} & \ar{دائرة المختلف} (1) \\
\textbf{Aruz (عروض)} & \ar{مقبوضة} (Qabadh: 5th element quiescent) → \ar{فَعُولُ} \\
\textbf{Dharb (ضرب)} & \ar{مقبوض}: \ar{فَعُولُ} \quad\textbar\quad \ar{محذوف}: \ar{فَعَلْ} \quad\textbar\quad \ar{صحيح}: \ar{مَفَاعِيلُنْ} \\
\textbf{Main Zehafs} & \ar{القبض}: removes 5th element of \ar{فَعُولُنْ} → \ar{فَعُولُ}\newline \ar{الثلم}: removes 1st element → \ar{مَفَاعِيلُنْ} → \ar{فَاعِيلُنْ} (rare) \\
\textbf{Frequency} & Most common meter in classical Arabic poetry (\textgreater40\% of Diwan) \\
\textbf{Typical Themes} & All purposes — this is the ``default'' meter for serious poetry \\
\bottomrule
\end{tabularx}
\end{table}

\begin{verseframe}
\centering
\ar{قِفَا نَبْكِ مِنْ ذِكْرَى حَبِيبٍ وَمَنْزِلِ \quad بِسِقْطِ اللِّوَى بَيْنَ الدَّخُولِ فَحَوْمَلِ}\\[4pt]
\small\textit{Imru' al-Qays, Mu'allaqah (opening verse)}
\end{verseframe}

% -----------------------------------------------------------------------
\subsection{\ar{البحر البسيط} — Bahr al-Baseet (``The Simple'')}
\label{subsec:baseet}

\begin{table}[H]
\centering
\caption{Al-Baseet structure}
\renewcommand{\arraystretch}{1.4}
\begin{tabularx}{\textwidth}{lX}
\toprule
\textbf{Hemistich} & \ar{مُسْتَفْعِلُنْ فَاعِلُنْ مُسْتَفْعِلُنْ فَاعِلُنْ} \\
\textbf{Pattern} & 1\,0\,1\,0\,1\,1\,0 \; 1\,0\,1\,1\,0 \; 1\,0\,1\,0\,1\,1\,0 \; 1\,0\,1\,1\,0 \\
\textbf{Circle} & \ar{دائرة المختلف} (1) \\
\textbf{Aruz} & \ar{مخبونة}: \ar{مُفَاعِلُنْ} (Khaban on \ar{مُسْتَفْعِلُنْ}) \\
\textbf{Dharb} & \ar{مخبون}: \ar{فَعِلُنْ} \quad\textbar\quad \ar{مقطوع}: \ar{فَعْلُنْ} \\
\textbf{Main Zehafs} & \ar{الخبن}: removes 2nd element (\ar{مُسْتَفْعِلُنْ} → \ar{مُفَاعِلُنْ})\newline \ar{الطي}: removes 4th element (\ar{مُسْتَفْعِلُنْ} → \ar{مُسْتَعِلُنْ}) \\
\textbf{Frequency} & Second most common meter \\
\textbf{Typical Themes} & Lament (\ar{رثاء}), description (\ar{وصف}), wisdom (\ar{حكمة}) \\
\bottomrule
\end{tabularx}
\end{table}

\begin{verseframe}
\centering
\ar{أَنَا الَّذِي نَظَرَ الأَعْمَى إِلَى أَدَبِي \quad وَأَسْمَعَتْ كَلِمَاتِي مَنْ بِهِ صَمَمُ}\\[4pt]
\small\textit{Al-Mutanabbi}
\end{verseframe}

% -----------------------------------------------------------------------
\subsection{\ar{البحر الكامل} — Bahr al-Kamel (``The Complete'')}
\label{subsec:kamel}

\begin{table}[H]
\centering
\caption{Al-Kamel structure}
\renewcommand{\arraystretch}{1.4}
\begin{tabularx}{\textwidth}{lX}
\toprule
\textbf{Hemistich} & \ar{مُتَفَاعِلُنْ مُتَفَاعِلُنْ مُتَفَاعِلُنْ} \\
\textbf{Pattern} & 1\,1\,1\,0\,1\,1\,0 \; 1\,1\,1\,0\,1\,1\,0 \; 1\,1\,1\,0\,1\,1\,0 \\
\textbf{Circle} & \ar{دائرة المؤتلف} (2) \\
\textbf{Aruz} & \ar{صحيحة}: \ar{مُتَفَاعِلُنْ} \quad\textbar\quad \ar{مُضمَرة}: \ar{مُتْفَاعِلُنْ} \\
\textbf{Dharb} & \ar{صحيح}, \ar{مُضمَر}, \ar{مقطوع}: \ar{مُتَفَاعِلْ}, \ar{أَحَث}: \ar{فَعِلَاتُنْ} \\
\textbf{Main Zehafs} & \ar{الإضمار}: vowels 2nd element (\ar{مُتَفَاعِلُنْ} → \ar{مُتْفَاعِلُنْ})\newline \ar{الوقص}: removes 2nd element (rare, only in hashw)\newline \ar{الحث}: combines إضمار + وقص on fasila \\
\textbf{Frequency} & Third most common; the meter of al-Mutanabbi's great odes \\
\textbf{Typical Themes} & Panegyric (\ar{مديح}), martial (\ar{حماسة}), boasting (\ar{فخر}) \\
\bottomrule
\end{tabularx}
\end{table}

\begin{verseframe}
\centering
\ar{وَعَلَى الرِّيَاحِ يَرُودُ نَبْلِي أَنَّنِي \quad قَوَّمْتُ مِنْ عِوَجِ الرَّدَى سَهْمًا مُقِيمَا}\\[4pt]
\small\textit{Al-Mutanabbi, from his Kafur odes}
\end{verseframe}

% -----------------------------------------------------------------------
\subsection{\ar{البحر الوافر} — Bahr al-Wafir (``The Abundant'')}
\label{subsec:wafir}

\begin{table}[H]
\centering
\caption{Al-Wafir structure}
\renewcommand{\arraystretch}{1.4}
\begin{tabularx}{\textwidth}{lX}
\toprule
\textbf{Hemistich} & \ar{مُفَاعَلَتُنْ مُفَاعَلَتُنْ مُفَاعَلَتُنْ} (full) \\
\textbf{Pattern} & 1\,1\,0\,1\,1\,1\,0 \; 1\,1\,0\,1\,1\,1\,0 \; (aruz variant) \\
\textbf{Circle} & \ar{دائرة المؤتلف} (2) \\
\textbf{Aruz/Dharb} & Always \ar{عصب} (\ar{مَقْطُوفة}): 5th + 6th elements quiescent → \ar{مُفَاعَلْتُنْ} → \ar{مَفَاعِيلُنْ} \\
\textbf{Main Zehafs} & \ar{العصب}: suppresses 5th element in \ar{مُفَاعَلَتُنْ} → \ar{مُفَاعِلْتُنْ} (mandatory in aruz/dharb) \\
\textbf{Note} & The \ar{عصب} modification makes the full form of \ar{مُفَاعَلَتُنْ} extremely rare in the last foot \\
\textbf{Frequency} & Common; associated with emotional and lyric poetry \\
\textbf{Typical Themes} & Love (\ar{غزل}), longing (\ar{شوق}), elegy (\ar{رثاء}) \\
\bottomrule
\end{tabularx}
\end{table}

% -----------------------------------------------------------------------
\subsection{\ar{البحر الرجز} — Bahr al-Rajaz (``The Trembling'')}
\label{subsec:rajaz}

\begin{table}[H]
\centering
\caption{Al-Rajaz structure}
\renewcommand{\arraystretch}{1.4}
\begin{tabularx}{\textwidth}{lX}
\toprule
\textbf{Hemistich} & \ar{مُسْتَفْعِلُنْ مُسْتَفْعِلُنْ مُسْتَفْعِلُنْ} \\
\textbf{Pattern} & 1\,0\,1\,0\,1\,1\,0 \; (× 3 per hemistich) \\
\textbf{Circle} & \ar{دائرة المجتلب} (3) \\
\textbf{Special Forms} & \ar{مجزوء}: 2 feet per hemistich \quad\textbar\quad \ar{مشطور}: 3 feet total (single shatr) \quad\textbar\quad \ar{منهوك}: 2 feet total \\
\textbf{Unique Feature} & The only meter in which a poet may compose alone (using \ar{مشطور}/\ar{منهوك} forms) without a response partner; this makes it the meter of \ar{الأراجيز} (didactic solo verse) \\
\textbf{Main Zehafs} & \ar{الخبن}, \ar{الطي}, \ar{الخبل} (= خبن + طي simultaneously) \\
\textbf{Frequency} & Very common; the original meter of \ar{حداء} (camel-driver's song) \\
\textbf{Typical Themes} & Didactic poetry (\ar{الأراجيز التعليمية}), warrior/martial, improvised verse \\
\bottomrule
\end{tabularx}
\end{table}

% -----------------------------------------------------------------------
\subsection{\ar{البحر الرمل} — Bahr al-Ramal (``The Sand'')}
\label{subsec:ramal}

\begin{table}[H]
\centering
\caption{Al-Ramal structure}
\renewcommand{\arraystretch}{1.4}
\begin{tabularx}{\textwidth}{lX}
\toprule
\textbf{Hemistich} & \ar{فَاعِلَاتُنْ فَاعِلَاتُنْ فَاعِلَاتُنْ} \\
\textbf{Pattern} & 1\,0\,1\,1\,0\,1\,0 \; (× 3 per hemistich) \\
\textbf{Circle} & \ar{دائرة المجتلب} (3) \\
\textbf{Aruz} & Always truncated (\ar{محذوف}): \ar{فَاعِلَاتُنْ} → \ar{فَاعِلَا} (= \ar{فَاعِلَاتٌ} with last element dropped) \\
\textbf{Dharb} & \ar{محذوف}: \ar{فَاعِلَا} \quad\textbar\quad \ar{مخبون}: \ar{فَعِلَاتُنْ} and variants \\
\textbf{Main Zehafs} & \ar{الخبن}: removes 2nd element \quad\textbar\quad \ar{الكف}: removes 7th element \\
\textbf{Frequency} & Common, especially in Andalusian and post-classical poetry \\
\textbf{Typical Themes} & Lyric poetry (\ar{غزل}), description (\ar{وصف}), Sufi poetry \\
\bottomrule
\end{tabularx}
\end{table}

% -----------------------------------------------------------------------
\subsection{\ar{البحر الهزج} — Bahr al-Hazaj (``The Buzzing'')}
\label{subsec:hazaj}

\begin{table}[H]
\centering
\caption{Al-Hazaj structure}
\renewcommand{\arraystretch}{1.4}
\begin{tabularx}{\textwidth}{lX}
\toprule
\textbf{Hemistich} & \ar{مَفَاعِيلُنْ مَفَاعِيلُنْ} \\
\textbf{Pattern} & 1\,1\,0\,1\,1\,0 \; 1\,1\,0\,1\,1\,0 \\
\textbf{Circle} & \ar{دائرة المجتلب} (3) \\
\textbf{Aruz} & Intact (\ar{سالمة}) or with kaf (\ar{مكفوفة}): \ar{مَفَاعِيلُ} \\
\textbf{Dharb} & \ar{سالم}: \ar{مَفَاعِيلُنْ} \quad\textbar\quad \ar{محذوف}: \ar{مَفَاعِيلْ} \\
\textbf{Main Zehafs} & \ar{الكف}: removes 7th element (position 6 of \ar{مَفَاعِيلُنْ}) \quad\textbar\quad \ar{القبض} disallowed in hashw \\
\textbf{Note} & The simplest meter — only two feet per hemistich; often used for songs \\
\textbf{Typical Themes} & Light songs (\ar{أغاني}), love (\ar{غزل}), folk verse \\
\bottomrule
\end{tabularx}
\end{table}

% -----------------------------------------------------------------------
\subsection{\ar{البحر الخفيف} — Bahr al-Khafeef (``The Light'')}
\label{subsec:khafeef}

\begin{table}[H]
\centering
\caption{Al-Khafeef structure}
\renewcommand{\arraystretch}{1.4}
\begin{tabularx}{\textwidth}{lX}
\toprule
\textbf{Hemistich} & \ar{فَاعِلَاتُنْ مُسْتَفْعِ لُنْ فَاعِلَاتُنْ} \\
\textbf{Pattern} & 1\,0\,1\,1\,0\,1\,0 \; 1\,0\,1\,0\,1\,0 \; 1\,0\,1\,1\,0\,1\,0 \\
\textbf{Circle} & \ar{دائرة المشتبه} (4) — contains watad mafruq \\
\textbf{Aruz} & Intact (\ar{سالمة}) or khabun: \ar{فَعِلَاتُنْ} \\
\textbf{Dharb} & \ar{سالم} \quad\textbar\quad \ar{مشعوث}: \ar{فَاعِلَاتُ} \quad\textbar\quad \ar{محذوف}: \ar{فَاعِلَا} \\
\textbf{Main Zehafs} & \ar{الخبن} on \ar{فَاعِلَاتُنْ} (positions 1,2) \\
\textbf{Note} & The watad mafruq (1\,0\,1) in \ar{مُسْتَفْعِ لُنْ} is a fixed constraint — its 2nd element (position 2 of the foot = position 2 of the bayt-middle foot) must remain quiescent \\
\textbf{Typical Themes} & Social/occasional verse, wisdom, religious poetry \\
\bottomrule
\end{tabularx}
\end{table}

% -----------------------------------------------------------------------
\subsection{\ar{البحر السريع} — Bahr al-Saree' (``The Fast'')}
\label{subsec:saree}

\begin{table}[H]
\centering
\caption{Al-Saree' structure}
\renewcommand{\arraystretch}{1.4}
\begin{tabularx}{\textwidth}{lX}
\toprule
\textbf{Hemistich} & \ar{مُسْتَفْعِلُنْ مُسْتَفْعِلُنْ مَفْعُولَاتُ} \\
\textbf{Pattern} & 1\,0\,1\,0\,1\,1\,0 \; 1\,0\,1\,0\,1\,1\,0 \; 1\,0\,1\,0\,1\,1\,0 \\
\textbf{Circle} & \ar{دائرة المشتبه} (4) \\
\textbf{Note on mafoolat} & \ar{مَفْعُولَاتُ} contains a watad mafruq; this foot only appears at the end \\
\textbf{Aruz/Dharb} & \ar{مَكْشُوفة/مَوْقُوفة}: \ar{مَفْعُولاتُ} takes \ar{وقف} (position 6 quiescent) → \ar{مَفْعُولَاتْ}; or \ar{كسف} (removes last element) → \ar{مَفْعُولَا} \\
\textbf{Main Zehafs} & \ar{الطي} on \ar{مُسْتَفْعِلُنْ} (removes 4th element) \quad\textbar\quad \ar{الخبل} (= خبن + طي) \\
\textbf{Typical Themes} & Heroic, panegyric \\
\bottomrule
\end{tabularx}
\end{table}

% -----------------------------------------------------------------------
\subsection{\ar{البحر المتقارب} — Bahr al-Mutakareb (``The Approaching'')}
\label{subsec:mutakareb}

\begin{table}[H]
\centering
\caption{Al-Mutakareb structure}
\renewcommand{\arraystretch}{1.4}
\begin{tabularx}{\textwidth}{lX}
\toprule
\textbf{Hemistich} & \ar{فَعُولُنْ فَعُولُنْ فَعُولُنْ فَعُولُنْ} \\
\textbf{Pattern} & 1\,1\,0\,1\,0 \; (× 4 per hemistich) \\
\textbf{Circle} & \ar{دائرة المتفق} (5) \\
\textbf{Aruz} & \ar{مقبوضة}: \ar{فَعُولُ} \quad\textbar\quad \ar{محذوفة}: \ar{فَعُلْ} (= \ar{فَعَلْ}) \\
\textbf{Dharb} & \ar{مقبوض}, \ar{محذوف}, \ar{أبتر}: \ar{فَعْلُنْ} \\
\textbf{Main Zehafs} & \ar{القبض}: removes 5th element → \ar{فَعُولُ}\newline \ar{الثلم}: removes 1st element (rare)\newline \ar{الثرم}: removes 1st + 2nd (very rare) \\
\textbf{Frequency} & Common in epic-narrative poetry; the meter of Ferdowsi's \textit{Shahnameh} equivalent in Arabic \\
\textbf{Typical Themes} & Epics, heroic narrative, occasional \\
\bottomrule
\end{tabularx}
\end{table}

\begin{verseframe}
\centering
\ar{أَلَا لَيْتَ الشَّبَابَ يَعُودُ يَوْمًا \quad فَأُخْبِرَهُ بِمَا فَعَلَ الْمَشِيبُ}\\[4pt]
\small\textit{Classical verse attributed to multiple poets}
\end{verseframe}

% -----------------------------------------------------------------------
\subsection{\ar{البحر المتدارك / المحدث} — Bahr al-Mutadarak}
\label{subsec:mutadarak}

\begin{table}[H]
\centering
\caption{Al-Mutadarak structure}
\renewcommand{\arraystretch}{1.4}
\begin{tabularx}{\textwidth}{lX}
\toprule
\textbf{Hemistich} & \ar{فَاعِلُنْ فَاعِلُنْ فَاعِلُنْ فَاعِلُنْ} \\
\textbf{Pattern} & 1\,0\,1\,1\,0 \; (× 4 per hemistich) \\
\textbf{Circle} & \ar{دائرة المتفق} (5) \\
\textbf{Historical Note} & Added by al-Akhfash al-Awsat after al-Khalil's death, hence called \ar{المحدث} (the newly added); al-Khalil recognised it but didn't include it \\
\textbf{Zehafs} & \ar{الخبن}: removes 2nd element → \ar{فَعِلُنْ} \\
\textbf{Typical Themes} & Songs, light poetry, modern verse \\
\bottomrule
\end{tabularx}
\end{table}

\noindent The remaining six meters — \ar{المديد}, \ar{المنسرح}, \ar{المضارع}, \ar{المقتضب}, \ar{المجتث} — are summarised in Appendix~\ref{app:meters_reference} with full structural data. They are rare in the corpus and primarily appear in specific thematic contexts (mujtath: occasional verse; mudhare: very rare; maqtadheb: nearly extinct in practice).

% ============================================================================
\section{The Rhyme System (\ar{نظام القافية})}
\label{sec:rhyme_system}

\subsection{The Five Rhyme Letters}
\label{subsec:five_rhyme_letters}

Arabic rhyme is not simply end-rhyme. A classical qafiya is defined by up to five letters that form a structural unit at the end of each verse:

\begin{table}[H]
\centering
\caption{The five qafiya letters (\ar{حروف القافية})}
\label{tab:qafiya_letters}
\renewcommand{\arraystretch}{1.5}
\begin{tabularx}{\textwidth}{lllX}
\toprule
\textbf{Letter} & \textbf{Arabic} & \textbf{Position} & \textbf{Definition} \\
\midrule
Al-Rawiy & \ar{الرَّوِيّ} & Mandatory & The main rhyme consonant; defines the rhyme entirely \\
Al-Wasl & \ar{الوَصْل} & Optional & Follows the rawiy; can be \ar{هاء}, long vowel, or n. of tanwin \\
Al-Khuruuj & \ar{الخُرُوج} & Optional & Vowel on the wasl consonant \\
Al-Radf & \ar{الرِّدْف} & Optional & Precedes rawiy; must be a long vowel (\ar{ا، و، ي}) or hamza \\
Al-Ta'siis & \ar{التَّأْسِيس} & Optional & \ar{أَلِف} two positions before rawiy, always with a consonant between \\
\bottomrule
\end{tabularx}
\end{table}

\subsection{Rhyme Vowel Types (\ar{حركات القافية})}

\begin{description}
  \item[\ar{المطلقة}] The rawiy bears a short vowel; the verse ends with an open syllable. Most common and musically richest.
  \item[\ar{المُقَيَّدة}] The rawiy bears a sukun; the verse ends with a closed syllable. More austere.
  \item[\ar{المُجَرَّدة}] The rawiy is bare (no radf, no ta'siis).
  \item[\ar{المُرَدَّفة}] The rawiy is preceded by a radf (long vowel before the rhyme consonant).
  \item[\ar{المُؤَسَّسة}] The rawiy is preceded by a ta'siis (alif two positions back).
\end{description}

\subsection{Qafiya Defects (\ar{عيوب القافية})}

\begin{table}[H]
\centering
\caption{Classical qafiya defects}
\label{tab:qafiya_defects}
\renewcommand{\arraystretch}{1.5}
\begin{tabularx}{\textwidth}{llX}
\toprule
\textbf{Name} & \textbf{Arabic} & \textbf{Definition} \\
\midrule
Al-'Ita' & \ar{الإيطاء} & Repeating the exact same word as rawiy in two or more verses (the worst defect) \\
Al-Ikfa' & \ar{الإكفاء} & Using two different letters that are similar in articulation as rawiy (\ar{ن/م} etc.) \\
Al-Sinnad & \ar{السِّناد} & Inconsistency in the radf or ta'siis between verses \\
Al-Iqwa' & \ar{الإقواء} & Mixing kasra and damma on the rawiy in the same poem \\
Al-'Israf & \ar{الإصراف} & Mixing fatha with kasra or damma on the rawiy \\
Al-Ijazah & \ar{الإجازة} & Changing the rawiy letter mid-poem \\
\bottomrule
\end{tabularx}
\end{table}

\subsection{Qafiya Type by Inter-Consonant Vowel Count}

\begin{table}[H]
\centering
\caption{Qafiya types by vowel structure}
\label{tab:qafiya_types}
\renewcommand{\arraystretch}{1.5}
\begin{tabularx}{\textwidth}{lllXl}
\toprule
\textbf{Type} & \textbf{Arabic} & \textbf{Vowels Between} & \textbf{Description} & \textbf{Example Ending} \\
\midrule
Mutawatir & \ar{مُتَوَاتِر} & One & One mutaharrik between the two saakins (long vowels) & \ar{شَجَنْ}، \ar{وَطَنْ}، \ar{عَزِيزَانِ}، \ar{حَبِيبَانِ} \\
Mutadaarik & \ar{مُتَدَارِك} & Two & Two mutaharrikaat between the two saakins & \ar{الْعَنْبَرِ}، \ar{قَمَرٌ}، \ar{كِتَابَةٌ} \\
Mutarakib & \ar{مُتَرَاكِب} & Three & Three mutaharrikaat between the two saakins & \ar{الْفَرَجِ}، \ar{الْحَسَبِ}، \ar{الْأَدَبِ} \\
Mutakaasis & \ar{مُتَكَاوِس} & Four & Four mutaharrikaat between the two saakins — very rare & \ar{فَجُبِرْ} \\
Mutaradif & \ar{مُتَرَادِف} & Zero & Two adjacent sukuns at the end — rawiy is maqayyada & \ar{دِينْ}، \ar{عَيْنْ}، \ar{جَوَادْ} \\
\bottomrule
\end{tabularx}
\end{table}

% ============================================================================
\section{Classical Qasida Structure (\ar{بنية القصيدة})}
\label{sec:qasida_structure}

The classical Arabic qasida (\ar{القصيدة الكلاسيكية}) follows a conventional tripartite or quadripartite structure, though not every poem adheres to it:

\begin{enumerate}
  \item \textbf{\ar{النسيب / المطلع}} — \textit{The amatory prelude / opening.} The poem traditionally opens with a lament at a deserted campsite (\ar{الوقوف على الأطلال}) or an evocation of the beloved's memory. This establishes an emotional key. The quality of the \ar{مطلع} (opening verse) was considered the most critical measure of a poem's success — it must arrest the listener immediately.

  \item \textbf{\ar{الرحيل / التخلص}} — \textit{The journey / transition.} The poet shifts from the opening mood toward the poem's main purpose, often through the image of a desert journey or a description of the poet's camel/horse.

  \item \textbf{\ar{الغرض الرئيسي}} — \textit{The main purpose section.} Panegyric, elegy, boasting, satire — this is where the declared purpose of the poem is fulfilled.

  \item \textbf{\ar{الخاتمة}} — \textit{The conclusion.} Often a prayer, a maxim, or a return to the opening image.
\end{enumerate}

\begin{rulebox}{Rules for a \ar{مطلع حسن} (Excellent Opening Verse)}
\begin{itemize}
  \item The opening verse must be complete in meaning — it must not require the next verse to make sense.
  \item It should establish the emotional tone of the entire poem.
  \item It should exhibit \ar{براعة الاستهلال} (masterly beginning): an image, sound, or idea that captures the listener's full attention immediately.
  \item It should avoid \ar{التعقيد اللفظي} (verbal complexity) — the opening must be immediately accessible.
  \item The first hemistich (\ar{الصدر}) establishes the rhyme letter (\ar{الروي}) for the entire poem. The second hemistich (\ar{العجز}) also rhymes with it — the first \ar{بيت} is \ar{مُصَرَّع} (both hemistichs rhyme).
\end{itemize}
\end{rulebox}

% ============================================================================
\section{Poetic Purposes (\ar{أغراض الشعر})}
\label{sec:aghard}

\subsection{\ar{المديح} — Panegyric}
\label{subsec:madih}

The most commercially and politically significant غرض in the classical tradition. A panegyric poet praised rulers, tribal chiefs, or patrons in exchange for gifts and protection.

\textbf{Structural Progression:} The classical \ar{مديحة} typically moves through: (1) an amatory or descriptive opening (\ar{نسيب}); (2) a transition (\ar{تخلص}) introducing the praised person; (3) the panegyric body enumerating the \textbf{four cardinal virtues}: generosity (\ar{الكرم}), valor (\ar{الشجاعة}), wisdom (\ar{الحلم}), justice (\ar{العدل}); (4) a supplicatory close.

\textbf{Language \& Style:} Grand, elevated diction; extensive use of superlative forms; comparison with natural phenomena (the praised person's generosity = the sea; courage = a lion or eagle); hyperbole (\ar{المبالغة}) is not a defect but a required feature.

\textbf{Meters Preferred:} al-Taweel, al-Baseet, al-Kamel, al-Wafir.

\textbf{Key Poets:} al-Mutanabbi (Abu al-Tayyib Ahmad ibn al-Husayn), al-Buhturi, Abu Tammam, al-Farazdaq, Hassan ibn Thabit.

\begin{verseframe}
\centering\ar{لَكَ يَا مَنَازِلُ فِي الْقُلُوبِ مَنَازِلُ \quad أَقْفَرْتِ أَنْتِ وَهُنَّ مِنْكِ أَوَاهِلُ}\\[4pt]
\small\textit{Al-Mutanabbi (Kamel meter)}
\end{verseframe}

\textbf{AI Generation Rules:}
\begin{itemize}
  \item The system must know \textit{who} is being praised and \textit{for what qualities} before generating.
  \item The standard three-phase structure (nasib → takhallus → madih body) must be maintained unless the user specifies otherwise.
  \item Hyperbole is correct in this غرض; the evaluator must not penalise it.
\end{itemize}

% ---
\subsection{\ar{الهجاء} — Satire/Invective}
\label{subsec:hija}

The inverse of madih: systematic attack on an enemy's character, lineage, appearance, and moral failings. Among the most feared forms of speech in early Arab culture — powerful enough that certain pre-Islamic practices surrounded its delivery with protective ritual.

\textbf{Sub-type — \ar{النقائض} (Poetic Flyting):} Two poets attack each other using the \textit{same meter and rhyme}, each poem a direct response to the other's. The great \ar{نقائض} exchanges of Jarir, al-Farazdaq, and al-Akhtal (c.700 CE) are the canonical examples. Rules: the responding poem \textit{must} use the same meter and rhyme as the poem being answered; the responding poet \textit{must} address every accusation before adding new ones.

\textbf{AI Generation Rules:}
\begin{itemize}
  \item For \ar{هجاء} as a new composition: the target's failings (stinginess = \ar{البخل}, cowardice = \ar{الجبن}, low birth = \ar{لؤم الأصل}) must be specified.
  \item For \ar{نقائض}: the system must read the original attack poem in full before generating the response; meter and rhyme are non-negotiable.
\end{itemize}

% ---
\subsection{\ar{الغزل} — Love Poetry}
\label{subsec:ghazal}

Divided into two fundamentally different schools:

\begin{description}
  \item[\ar{الغزل العذري}] \textit{Udhri/Platonic love}: Chaste, idealised, often unrequited love. Named after the Banu Udhrah tribe of the Hijaz. Characterised by: absolute prohibition on describing the beloved's body; the lover as a willing martyr to love; the beloved as unattainable ideal. Canonical figures: Qays and Layla (\ar{مجنون ليلى}), Jamil and Buthayna (\ar{جميل بثينة}), Kuthayyir and Azza.
  \item[\ar{الغزل الحسي / الصريح}] \textit{Sensual/explicit love}: Bold physical description; celebration of desire and consummation. The Hijazi/Umayyad school centred on \textbf{Umar ibn Abi Rabi'a} (d.~712 CE). Meters: Ramal, Sari', Hazaj preferred.
\end{description}

\textbf{AI Generation Rules:}
\begin{itemize}
  \item The system must distinguish between \ar{عذري} and \ar{حسي} modes — they have different constraints on body imagery.
  \item In \ar{عذري}: body description is prohibited; emotional and psychological states are the content.
  \item In \ar{حسي}: physical description is not only permitted but expected.
  \item The meter itself often signals the mode: Taweel/Baseet → more likely \ar{عذري}; Ramal/Hazaj → more likely \ar{حسي}.
\end{itemize}

% ---
\subsection{\ar{الرثاء} — Elegy}
\label{subsec:ritha}

Mourning poetry for the dead. Structurally, \ar{رثاء} is praise of the deceased combined with grief at the loss — it shares vocabulary and technique with \ar{مديح} but in a minor emotional key.

\textbf{Key feature:} The best \ar{رثاء} alternates between \ar{الحزن} (grief) and \ar{التمجيد} (glorification). Pure lament without glorification is a defect.

\textbf{Famous poetess:} \textbf{al-Khansa'} (\ar{الخنساء}) — her elegies for her brothers Sakhr and Mu'awiya remain the canonical model.

\textbf{AI Generation Rules:}
\begin{itemize}
  \item The AI must know who is being mourned (person, relationship, qualities to glorify) and whether the tone should be personal grief or communal mourning.
  \item Meter: al-Taweel and al-Baseet are traditional; al-Wafir is also common.
\end{itemize}

% ---
\subsection{\ar{الفخر} — Boasting/Self-Glorification}
\label{subsec:fakhr}

Glorification of the self, one's tribe, and one's deeds. The dominant tone is elevated pride, not mere self-promotion. In pre-Islamic and Umayyad poetry, personal \ar{فخر} was inseparable from tribal identity.

\subsection{\ar{الحكمة} — Wisdom Poetry}
\label{subsec:hikma}

Philosophical and moral aphorisms in verse. Brief, memorable, self-contained verses. Famous: Zuhair ibn Abi Sulma's \ar{معلقة}, Imru' al-Qays's philosophical lines, the gnomic verses of Labid.

\textbf{AI generation note:} \ar{حكمة} verses are typically stand-alone — a two-hemistich \ar{بيت} that delivers a complete insight. The AI should generate complete, pithy, universal statements.

\subsection{\ar{الزهد} — Asceticism}
\label{subsec:zuhd}

Renunciation of worldly pleasures and focus on the afterlife. Associated with: Abu al-Atahiya (d.~828 CE), the most prolific \ar{زهد} poet. Language is plain, direct, and rhythmically simple — deliberately avoiding the grandeur of \ar{مديح}.

\subsection{\ar{الخمريات} — Wine Poetry}
\label{subsec:khamriyyat}

Celebration of wine and drinking culture. Abu Nuwas (d.~c.815 CE) is the supreme master. \textbf{Note for AI:} In Sufi poetry (Ibn al-Farid, Omar Khayyam in Arabic translation, etc.), wine lexicon (\ar{خمر}, \ar{كأس}, \ar{ساقي}) carries a spiritual allegorical meaning (wine = divine knowledge; the cup-bearer = the Sufi master). The AI must distinguish literal from allegorical register.

\subsection{\ar{الشعر الصوفي} — Sufi/Mystical Poetry}
\label{subsec:sufi}

A major tradition from the 9th century onward. Key figures: al-Hallaj, Rabi'a al-Adawiyya, Ibn Arabi, Ibn al-Farid. Characteristics: deliberate ambiguity between human and divine love; lexical borrowing from \ar{غزل} and \ar{خمريات}; mystical paradox (\ar{شطح}).

\textbf{Core Sufi symbolic system (for AI evaluation):}
\begin{table}[H]
\centering
\caption{Sufi symbolic lexicon}
\label{tab:sufi_symbols}
\renewcommand{\arraystretch}{1.4}
\begin{tabularx}{\textwidth}{llX}
\toprule
\textbf{Literal Image} & \textbf{Arabic} & \textbf{Sufi Referent} \\
\midrule
Wine (\ar{الخمر}) & & Divine knowledge / gnosis \\
Tavern (\ar{الحانة}) & & The Sufi tekke / path \\
Cup-bearer (\ar{الساقي}) & & The Sufi master / God \\
Drunkenness (\ar{السُّكر}) & & Ecstatic union (\ar{وجد}) \\
The Beloved (\ar{الحبيب / المحبوب}) & & The Divine \\
\bottomrule
\end{tabularx}
\end{table}

% ============================================================================
\section{Intertextual Compositional Forms (\ar{الأشكال التناصية})}
\label{sec:intertextual}

These forms are compositional operations that relate a new poem to an existing one. They are among the most technically demanding requests the system will receive.

\subsection{\ar{المعارضة} — Response/Counter-Poem}
\label{subsec:muarada}

A new poem composed using the \textbf{same meter and same rhyme letter} as an existing poem, with the implicit or explicit claim of matching or surpassing the original.

\begin{rulebox}{Rules for \ar{المعارضة}}
\begin{enumerate}
  \item \textbf{Meter}: Must be identical — the same Khalilian meter; even the sub-type (e.g., if original uses \ar{مجزوء البسيط}, the response must too).
  \item \textbf{Rhyme}: Must use the \textit{exact same rawiy letter}. The harakah on the rawiy (the vowel type) must also match.
  \item \textbf{Theme}: Can differ — the معارضة may take a completely different purpose (\ar{غرض}) from the original, or may engage with the same themes.
  \item \textbf{Length}: Need not match the original exactly, but a معارضة shorter than the original is considered inferior by convention.
  \item \textbf{Quality criterion}: The معارضة succeeds when it is judged to match or exceed the original's quality. This requires the AI to read the full source poem before generating.
\end{enumerate}
\end{rulebox}

\textbf{Famous examples:} Al-Buhturi's معارضة of al-Mutanabbi; neo-classical poets Ahmad Shawqi and Hafiz Ibrahim wrote معارضات of Imru' al-Qays's Mu'allaqah.

\subsection{\ar{المساجلة} — Poetic Competition}
\label{subsec:musajala}

Two or more poets exchange verses in turn, each attempting to outdo the other in wit, imagery, and prosodic mastery. A \ar{مساجلة} is not merely back-and-forth — it follows rules:

\begin{rulebox}{Rules for \ar{المساجلة}}
\begin{enumerate}
  \item Both poets use the same meter and rhyme throughout the exchange.
  \item Each poet's entry must respond to the specific \textit{argument} or \textit{image} of the previous poet — not simply continue the same subject.
  \item The exchange continues until one poet is unable to respond or the agreed number of turns is complete.
  \item A poet who changes the meter or rhyme mid-exchange concedes defeat.
\end{enumerate}
\end{rulebox}

\subsection{\ar{التضمين} — Intertextual Borrowing}
\label{subsec:tadmin}

Incorporating a full verse or hemistich from another poet into one's own poem, where it becomes structurally integrated into the new poem's flow.

\begin{description}
  \item[\ar{تضمين}] Borrowing from another poet's verse. May be acknowledged (\ar{تضمين صريح}) or left for the reader to recognise (\ar{تضمين خفي}).
  \item[\ar{اقتباس}] Borrowing from the Quran or Hadith. Treated differently — governed by additional literary and religious norms.
  \item[\ar{تلميح}] Mere allusion to a famous text without direct quotation — the weakest form of intertextuality.
\end{description}

\textbf{AI rule:} The borrowed hemistich must fit seamlessly — same meter, same rhyme position if it's at the end of a verse. The surrounding context must make the borrowing feel intentional and meaningful, not forced.

\subsection{\ar{التشطير} — Hemistichal Splitting}
\label{subsec:tashtir}

Taking an existing poem and inserting new hemistichs into each verse, effectively doubling the length of each \ar{بيت}. The original hemistichs are \textit{immutable}; the poet inserts a new hemistich before each original hemistich.

\textbf{Structure of a تشطير verse:}
\begin{verbatim}
  [New hemistich A] + [Original hemistich 1 (sadr)] +
  [New hemistich B] + [Original hemistich 2 (ajuz)]
\end{verbatim}

\begin{rulebox}{Rules for \ar{التشطير}}
\begin{enumerate}
  \item The original verses are reproduced without any modification.
  \item New hemistichs use the same meter as the original.
  \item New hemistichs may or may not rhyme with the original's rawiy — if they do, the effect is richer.
  \item The semantic flow from new hemistich → original hemistich must be smooth and logical.
  \item The inserted hemistichs must not contradict or undermine the meaning of the original.
\end{enumerate}
\end{rulebox}

\subsection{\ar{التخميس} — Quintuple Form}
\label{subsec:takhmiis}

Adding three new hemistichs to each original \ar{بيت}, creating a five-line unit (\ar{خُمَاسية}). The structure is: AAA-BC where A = three new lines sharing an internal rhyme, B = original sadr, C = original ajuz.

\begin{rulebox}{Rules for \ar{التخميس}}
\begin{enumerate}
  \item The three new lines (A, A, A) share a rhyme among themselves, different from the original poem's rhyme.
  \item The third new line must lead naturally into the original \ar{صدر} — it is the transition.
  \item The original \ar{صدر} (B) and \ar{عجز} (C) are reproduced exactly.
  \item The new lines use the same meter as the original.
\end{enumerate}
\end{rulebox}

\subsection{\ar{إكمال القصيدة} — Completing an Unfinished Poem}
\label{subsec:ikmal}

Completing a poem that the original poet left unfinished, or that exists only in fragmentary form.

\begin{rulebox}{Rules for \ar{إكمال القصيدة}}
\begin{enumerate}
  \item \textbf{Meter and rhyme} must be maintained without exception.
  \item \textbf{Emotional register and tone} must match the existing verses — if the opening is mournful, the completion cannot shift to celebration.
  \item \textbf{Sentence length patterns}: if the original uses long, flowing sentences, the completion should not suddenly shift to short staccato hemistichs.
  \item \textbf{Vocabulary level}: if the original uses archaic/Bedouin diction (\ar{لغة بدوية}), the completion must maintain this register.
  \item \textbf{Imagery density}: if the original is rich in simile and metaphor, the completion must match that density.
  \item The completion should provide a convincing \ar{خاتمة} — a sense of closure.
\end{enumerate}
\end{rulebox}

\subsection{\ar{التصوير الشعري} — Poetic Visualization}
\label{subsec:tasweer}

Translating a scene, experience, or concept into verse using the full range of Arabic rhetorical devices (\ar{البلاغة}).

\begin{table}[H]
\centering
\caption{Primary Arabic rhetorical devices (\ar{الأساليب البلاغية})}
\label{tab:rhetorical_devices}
\renewcommand{\arraystretch}{1.5}
\begin{tabularx}{\textwidth}{llX}
\toprule
\textbf{Device} & \textbf{Arabic} & \textbf{Definition and AI Generation Note} \\
\midrule
Simile & \ar{التشبيه} & Explicit comparison (A is like B); requires كأن, مثل, or equivalent. Most common device. \\
Metaphor & \ar{الاستعارة} & Implicit comparison; A is called B without the comparison particle. \\
Metonymy & \ar{الكناية} & Describing A by describing its characteristic attribute B (``heavy-handed'' for generous). \\
Synecdoche & \ar{المجاز المرسل} & Part for whole or whole for part; figurative extension. \\
Antithesis & \ar{الطباق} & Pairing contrasting concepts in the same verse. Highly valued in Abbasid poetry. \\
Paranomasia & \ar{الجناس} & Wordplay on similar-sounding words with different meanings. \\
\bottomrule
\end{tabularx}
\end{table}

\subsection{\ar{شعر المناسبات} — Occasional Poetry}
\label{subsec:munasabat}

Poetry composed for specific social occasions. Subtypes:
\begin{itemize}
  \item \ar{التهنئة} (Congratulatory): Weddings, births, promotions, achievements. Tone: joyful, elevated.
  \item \ar{العزاء / التعزية} (Condolence): Death, illness, misfortune. Tone: measured grief, comfort.
  \item \ar{المديح الرسمي} (Official praise): State occasions, national days. Tone: formal, civic.
  \item \ar{الوصف التصويري} (Descriptive occasion): Describing a place, event, season. Tone: varies.
\end{itemize}

% ============================================================================
\clearpage
\langsep{}

% ============================================================================
%  ARABIC VERSION
% ============================================================================
\begin{otherlanguage}{arabic}

\section*{الشعر العربي الفصيح: النظام العروضي والأغراض}

\subsection*{أسس علم العروض}

علم العروض هو العلم الذي وضعه الخليل بن أحمد الفراهيدي (118–170هـ / 736–786م) لضبط أوزان الشعر العربي. يقوم هذا العلم على ثلاثة مفاهيم أساسية:

\begin{description}
  \item[\textbf{السبب الخفيف}] حرفان: الأول متحرك والثاني ساكن. مثل: مَنْ، هَلْ، قَدْ.
  \item[\textbf{الوتد المجموع}] ثلاثة أحرف: متحرك + متحرك + ساكن. مثل: فَعُو، قَفَا.
  \item[\textbf{الوتد المفروق}] ثلاثة أحرف: متحرك + ساكن + متحرك. مثل: مَفْعُ.
\end{description}

تتركب من هذه المكونات ثمانُ تفعيلاتٍ أساسية، ومن هذه التفعيلات تتشكل ستة عشر بحرًا خليليًا.

\subsection*{البحور الخليلية الستة عشر}

فيما يلي ملخص البحور الأكثر استخدامًا مع أمثلتها:

\textbf{البحر الطويل} — أكثر البحور شيوعًا في الشعر العربي الكلاسيكي (أكثر من 40\% من الديوان الشعري). تفعيلاته: \textbf{فَعُولُنْ مَفَاعِيلُنْ فَعُولُنْ مَفَاعِيلُنْ} في كل شطر.

مثال: \textbf{قِفَا نَبْكِ مِنْ ذِكْرَى حَبِيبٍ وَمَنْزِلِ --- بِسِقْطِ اللِّوَى بَيْنَ الدَّخُولِ فَحَوْمَلِ}

(امرؤ القيس، مطلع المعلقة)

\textbf{البحر الكامل} — ثالث البحور شيوعًا، متفعيلاته: \textbf{مُتَفَاعِلُنْ مُتَفَاعِلُنْ مُتَفَاعِلُنْ}. المتنبي نظم أعظم قصائده على هذا البحر.

\textbf{البحر البسيط} — تفعيلاته: \textbf{مُسْتَفْعِلُنْ فَاعِلُنْ مُسْتَفْعِلُنْ فَاعِلُنْ}. يُستخدم كثيرًا في الرثاء والحكمة.

\subsection*{نظام القافية}

القافية في الشعر العربي الكلاسيكي ليست مجرد جناس صوتي في نهاية الشطر، بل هي منظومة دقيقة من الحروف والحركات:

\begin{description}
  \item[\textbf{الرَّوِيّ}] الحرف الأساسي للقافية؛ يحدد هوية القصيدة كلها.
  \item[\textbf{الوَصْل}] ما يلي الروي من هاء أو مدّ أو نون تنوين.
  \item[\textbf{الرِّدْف}] مدّ يسبق الروي مباشرة (ألف أو واو أو ياء).
  \item[\textbf{التَّأْسِيس}] ألف تفصلها عن الروي حرف واحد.
\end{description}

\textbf{عيوب القافية الكبرى:}
\begin{itemize}
  \item \textbf{الإيطاء}: تكرار كلمة الروي نفسها في أكثر من بيت (أقبح العيوب).
  \item \textbf{الإقواء}: الخلط بين الكسرة والضمة على الروي.
  \item \textbf{السناد}: التناقض في الردف أو التأسيس بين الأبيات.
\end{itemize}

\subsection*{أغراض الشعر العربي الكلاسيكي}

\textbf{المديح} — أعظم الأغراض وأوسعها في الشعر الكلاسيكي. يمر عبر ثلاث مراحل: النسيب (التمهيد)، ثم التخلص (الانتقال إلى الممدوح)، ثم جسم المديح (ذكر الكرم والشجاعة والحلم والعدل). قاعدة للذكاء الاصطناعي: يجب أن يعرف النظام بوضوح مَن يُمدَح ولأي صفات.

\textbf{الهجاء} — نقيض المديح؛ هجاء الأعداء وفضح عيوبهم. أهم أشكاله الخاصة: \textbf{النقائض} بين الشعراء؛ وفيها يجب أن يستخدم الشاعر المجيب نفس الوزن والقافية.

\textbf{الغزل العذري} — الحب العفيف المحتشم. يُحرَّم فيه وصف جسد المحبوب. نماذجه: مجنون ليلى، جميل بثينة.

\textbf{الغزل الحسي} — الحب الجريء الصريح. مدرسة عمر بن أبي ربيعة الحجازية. يُجيز وصف الجسد صراحة.

\textbf{الرثاء} — الحزن على الميت مع تمجيد مناقبه. الخنساء أعظم رثّاءة في تاريخ الشعر العربي.

\textbf{الزهد} — التبرؤ من الدنيا والتزهد فيها. أبو العتاهية سيد هذا الغرض.

\textbf{الخمريات} — الاحتفاء بالخمر والشرب. أبو نواس إمامها. تنبيه للذكاء الاصطناعي: في الشعر الصوفي، لغة الخمر رمزية روحانية لا حرفية.

\textbf{الشعر الصوفي} — الشعر الروحاني والعرفاني. الخمر = المعرفة الإلهية؛ الساقي = الشيخ أو الحق؛ السُّكر = حال الوجد.

\subsection*{الأشكال التناصية}

\textbf{المعارضة}: قصيدة جديدة على نفس وزن وقافية قصيدة أخرى. القاعدة الذهبية: الوزن والروي \textit{واجبان} لا خيار فيهما.

\textbf{التشطير}: أخذ قصيدة موجودة وإضافة شطر جديد قبل كل شطر من أشطارها الأصلية. الشروط: الشطور الأصلية لا تُعدَّل، والشطور الجديدة على نفس الوزن.

\textbf{التخميس}: إضافة ثلاثة أشطار جديدة لكل بيت أصلي، لتصبح الوحدة خمسة أشطار (أبب --- جج). الأشطار الثلاثة الجديدة تتقافى فيما بينها بقافية مختلفة عن قافية القصيدة الأصلية.

\textbf{إكمال القصيدة}: إتمام قصيدة ناقصة. الشروط: الوزن والقافية والنبرة العاطفية ومستوى اللغة والكثافة البلاغية كلها يجب أن تتطابق مع ما هو موجود.

\end{otherlanguage}

